\documentclass[aspectratio=169,10pt,t]{beamer}
\usepackage{ac1}
\aulainfo{21}{NumPy}{Vetores, operações vetorizadas e aplicações numéricas}

\begin{document}
\capa

\begin{frame}{Roteiro}
  \begin{itemize}
    \item Por que NumPy? Listas \emph{versus} arrays (Cap.~5.14)
    \item Criando arrays: \texttt{array}, \texttt{arange}, \texttt{linspace}, \texttt{zeros}, \texttt{ones}
    \item Operações vetorizadas e funções universais
    \item Indexação, fatiamento e filtros booleanos
    \item Matrizes, eixos (\texttt{axis}) e álgebra linear
    \item Armadilhas: homogeneidade e visões de memória
  \end{itemize}
\end{frame}

\begin{frame}{Por que NumPy?}
  \begin{columns}[T]
    \column{0.5\textwidth}
    \begin{block}{Listas do Python}
      \begin{itemize}\small
        \item Flexíveis (tipos misturados)
        \item \texttt{+} concatena, \texttt{*} repete
        \item Cálculos exigem \emph{laços} \(\Rightarrow\) lentos
      \end{itemize}
    \end{block}
    \column{0.5\textwidth}
    \begin{exampleblock}{Arrays do NumPy}
      \begin{itemize}\small
        \item Elementos de \textbf{um só tipo}
        \item Operações aplicadas elemento a elemento
        \item Código em C por baixo: \alert{muito} mais rápido
      \end{itemize}
    \end{exampleblock}
  \end{columns}
  \vspace{1ex}
  É a base da computação científica em Python: \textbf{Pandas}, \textbf{Matplotlib}, SciPy e \emph{machine learning} são construídos sobre NumPy.
\end{frame}

\begin{frame}[fragile]{Concatenar \emph{versus} somar}
\begin{lstlisting}
import numpy as np

a = [1, 2, 3]
b = [4, 5, 6]
c = a + b
print(c)               # [1, 2, 3, 4, 5, 6]   concatenação

a_num = np.array(a)
b_num = np.array(b)
c_num = a_num + b_num
print(c_num)           # [5 7 9]              soma vetorial
\end{lstlisting}
  \begin{itemize}
    \item \texttt{import numpy as np}: apelido universal
    \item Com arrays, \texttt{+}, \texttt{-}, \texttt{*}, \texttt{/}, \texttt{**} atuam \alert{elemento a elemento}
  \end{itemize}
  \fonte{code/Cap5/prog.5.h\_algebralinear.py}
\end{frame}

\begin{frame}[fragile]{Criando arrays}
\begin{lstlisting}
import numpy as np
a = np.array([1, 2, 3])                # de uma lista
print(np.arange(0, 1, 0.25))           # [0.   0.25 0.5  0.75]
print(np.linspace(0, 1, 5))            # [0.   0.25 0.5  0.75 1.  ]
print(np.zeros(3))                     # [0. 0. 0.]
print(np.ones((2, 3)))                 # matriz 2 x 3 de uns
print(np.eye(2))                       # identidade 2 x 2
print(a.shape, a.dtype, a.size)        # (3,) int64 3
\end{lstlisting}
  \begin{itemize}
    \item \texttt{arange(ini, fim, passo)}: como \texttt{range}, mas aceita passo \emph{real}
    \item \texttt{linspace(ini, fim, n)}: \(n\) pontos igualmente espaçados, \alert{incluindo} o fim
    \item \texttt{shape} = dimensões; \texttt{dtype} = tipo dos elementos
  \end{itemize}
\end{frame}

\begin{frame}[fragile]{Operações vetorizadas}
\begin{lstlisting}
a = np.array([1, 2, 3, 4])
print(a * 2)            # [2 4 6 8]
print(a ** 2)           # [ 1  4  9 16]
print(a + 10)           # [11 12 13 14]
print(1 / a)            # [1.  0.5  0.333...  0.25]
print(np.sqrt(a))       # [1.  1.41421356 1.73205081 2.]
x = np.linspace(0, np.pi, 5)
print(np.sin(x).round(3))    # [0.    0.707 1.    0.707 0.   ]
\end{lstlisting}
  \begin{itemize}
    \item Sem laço \texttt{for}: o código fica curto, claro e rápido
    \item \textbf{Funções universais} (\texttt{ufuncs}): \texttt{sin}, \texttt{cos}, \texttt{exp}, \texttt{log}, \texttt{sqrt}\ldots
    \item Uma tabela de senos (Aula 12) agora são duas linhas
  \end{itemize}
\end{frame}

\begin{frame}[fragile]{Velocidade: lista \emph{versus} array}
\begin{lstlisting}
import numpy as np, random, time
num = 10_000_000
lista1 = random.choices(range(1, 10001), k=num)
lista2 = random.choices(range(1, 10001), k=num)

t1 = time.time()
lista3 = [lista1[i] * lista2[i] for i in range(num)]
print('Lista:', time.time() - t1)

l1, l2 = np.array(lista1), np.array(lista2)
t1 = time.time()
l3 = l1 * l2
print('NumPy:', time.time() - t1)
\end{lstlisting}
  \small Em geral o NumPy é dezenas de vezes mais rápido (confira na sua máquina!).\\
  \fonte{code/Cap5/prog.5.i\_compara\_listas\_e\_arrays.py}
\end{frame}

\begin{frame}[fragile]{Estatística e agregações}
\begin{lstlisting}
notas = np.array([5, 9, 5.5, 10, 8])
print(notas.sum(), notas.mean())       # 37.5  7.5
print(notas.min(), notas.max())        # 5.0  10.0
print(notas.std(), np.median(notas))   # desvio-padrão, mediana
print(notas.argmax())                  # 3 (posição do maior)
print(np.sort(notas))                  # [ 5.   5.5  8.   9.  10. ]
\end{lstlisting}
  \textbf{Média ponderada} (álgebra linear em uma linha):
\begin{lstlisting}
pesos = np.array([1, 2, 4, 1, 2])
print(np.sum(notas * pesos) / np.sum(pesos))   # 7.1
print(np.average(notas, weights=pesos))        # 7.1
\end{lstlisting}
\end{frame}

\begin{frame}[fragile]{Indexação e filtros booleanos}
\begin{lstlisting}
import numpy as np
	
notas = np.array([5, 9, 5.5, 10, 8.4])
print(notas > 6)                   # [False  True False  True  True]
print('Notas > 6:', notas[notas > 6])      # [ 9.  10.   8.4]
print('Notas <= 6:', notas[notas <= 6])    # [5.  5.5]
print(((notas > 6) & (notas < 10)).sum())  # entre 6 e 10: 2
notas[notas < 5.5] = 5.5           # altera somente os que satisfazem
\end{lstlisting}
  \begin{itemize}
    \item Uma \textbf{máscara booleana} dentro dos colchetes seleciona elementos
    \item Combine com \texttt{\&} (e), \texttt{|} (ou), \texttt{\textasciitilde} (não), \alert{com parênteses}
    \item É a base de filtros em Pandas (Aula 28)
  \end{itemize}
  \fonte{code/Cap5/prog.5.g\_numpy\_turma.boolean.py}
\end{frame}

\begin{frame}[fragile]{Matrizes: \texttt{shape}, fatias e eixos}
\begin{lstlisting}
import numpy as np
	
m = np.array([[1, 2, 3],
              [4, 5, 6]])
print(m.shape)            # (2, 3)
print(m[1, 2])            # 6        (linha 1, coluna 2)
print(m[:, 0])            # [1 4]    (coluna 0)
print(m[0, :])            # [1 2 3]  (linha 0)
print(m.sum(axis=0))      # [5 7 9]  (soma de cada COLUNA)
print(m.sum(axis=1))      # [ 6 15]  (soma de cada LINHA)
print(m.T)                # transposta
\end{lstlisting}
  \begin{itemize}
    \item \texttt{m[i, j]} em vez de \texttt{m[i][j]}
    \item \texttt{axis=0} colapsa as \emph{linhas}; \texttt{axis=1}, as \emph{colunas}
  \end{itemize}
  \fonte{code/Cap5/prog.5.f\_numpy\_com\_axis.py}
\end{frame}

\begin{frame}[fragile]{Álgebra linear com NumPy}
  Resolver o sistema \(\begin{cases}2x + y = 3\\ x + 3y = 5\end{cases}\):
\begin{lstlisting}
A = np.array([[2, 1],
              [1, 3]])
b = np.array([3, 5])
x = np.linalg.solve(A, b)
print(x)                  # [0.8 1.4]
print(A @ x)              # [3. 5.]   (produto matricial; confere!)
print(np.linalg.det(A))   # 5.0 (aprox.)
print(np.linalg.inv(A))   # matriz inversa
\end{lstlisting}
  \begin{itemize}
    \item \texttt{*} é elemento a elemento; \texttt{@} é produto \textbf{matricial}
    \item Aplicações: circuitos (leis de Kirchhoff), estruturas, balanços de massa
  \end{itemize}
\end{frame}

\begin{frame}[fragile]{Aplicação: ajuste de reta (dados experimentais)}
\begin{lstlisting}
V = np.array([5.0, 10.0, 15.0, 20.0])     # tensão (V)
I = np.array([0.25, 0.52, 0.76, 1.01])    # corrente (A)

a, b = np.polyfit(V, I, 1)                # reta I = a V + b
print(f"I = {a:.4f} V + {b:.4f}")   # I = 0.0504 V + 0.0050
R = 1 / a
print(f"Resistência estimada: {R:.2f} ohms")   # 19.84
\end{lstlisting}
  \begin{itemize}
    \item A mesma tabela de ensaios da Aula 19, agora com análise numérica de verdade
    \item Na Aula 27 vamos \alert{plotar} os pontos e a reta com Matplotlib
  \end{itemize}
\end{frame}

\begin{frame}[fragile]{Armadilhas do NumPy}
  \textbf{1. Homogeneidade:} todos os elementos têm um tipo.
\begin{lstlisting}
print(np.array(['Python', 1, 2, 3]))   # ['Python' '1' '2' '3']  (tudo vira texto!)
\end{lstlisting}
  \textbf{2. Fatias são \emph{visões}, não cópias:}
\begin{lstlisting}
import numpy as np

a = [1, 2, 3, 4, 5];            b = a[1:];  b[0] = 42   # lista: a NÃO muda
a = np.array([1, 2, 3, 4, 5]);  b = a[1:];  b[0] = 42   # array: a MUDA!
print(a)                         # [ 1 42  3  4  5]
b = a[1:].copy()                 # para copiar de verdade
\end{lstlisting}
  \fonte{code/Cap5/prog.5.43\_numpy\_erros.py}
\end{frame}

\begin{frame}{Atividade}
  \begin{enumerate}
    \item Crie um array com 100 pontos de \(0\) a \(2\pi\) e calcule \(\sin\) e \(\cos\); verifique que \(\sin^2+\cos^2=1\) para todos os pontos.
    \item Dado o array de temperaturas de 24 horas, calcule média, máxima, hora da máxima, e quantas horas ficaram acima da média.
    \item Resolva o circuito de duas malhas (sistema \(3\times3\)) com \texttt{np.linalg.solve}.
    \item Reescreva, com NumPy, o programa Monte Carlo para estimar \(\pi\) sem nenhum laço \texttt{for}:\\
          \texttt{x = np.random.rand(n); y = np.random.rand(n); 4 * np.mean(x**2 + y**2 < 1)}
  \end{enumerate}
\end{frame}

\exercicios{5}{%
  \begin{itemize}\small
    \item Seção 5.14 (NumPy) do livro; as médias ponderadas e filtros da aula podem ser comparados com os exercícios de listas \textbf{5.7 a 5.13}
    \item Refaça com NumPy: \textbf{5.14} (fatiamento negativo), \textbf{5.25} (divisíveis por 7 e não por 3)
  \end{itemize}}

\resumo{
  \item \texttt{np.array}: elementos de um só tipo, operações elemento a elemento
  \item \texttt{arange}, \texttt{linspace}, \texttt{zeros}, \texttt{ones}; funções universais \texttt{sin}, \texttt{exp}, \texttt{sqrt}\ldots
  \item Filtros booleanos: \texttt{a[a > 6]}; \texttt{axis} define em que direção agregar
  \item Álgebra linear: \texttt{@}, \texttt{np.linalg.solve}, \texttt{polyfit}
  \item Cuidado: fatias são visões; use \texttt{.copy()}
  \item Próxima aula: strings
}

\end{document}
